\section*{Chapter \ref{ch:b3:inorganic-materials} --- Inorganic Materials and Nanomaterials}

\begin{solution}{exo:b3:inorganic-materials:1}
$x^2 \propto t$: three times thicker takes nine times longer, \qty{90}{h}.
\end{solution}

\begin{solution}{exo:b3:inorganic-materials:2}
Hydrolysis: $\ce{Si(OC2H5)4 + 4H2O -> Si(OH)4 + 4C2H5OH}$.\\
Condensation: $\ce{2Si(OH)4 -> Si2O(OH)6 + H2O}$.\\
Overall: $\ce{Si(OC2H5)4 + 2H2O -> SiO2 + 4C2H5OH}$.
\end{solution}

\begin{solution}{exo:b3:inorganic-materials:3}
Formers: $\ce{SiO2}$, $\ce{B2O3}$. Modifiers: $\ce{Na2O}$, $\ce{CaO}$. Each oxide ion of $\ce{Na2O}$ breaks one Si--O--Si
bridge into two Si--O$^-$ ends, each paired with a $\ce{Na+}$: the network is cut, the melt
flows more easily and melts lower.
\end{solution}

\begin{solution}{exo:b3:inorganic-materials:4}
Twelve pentagons (as for every \omterm{def:b3:inorganic-materials:carbon}{fullerene}). $V = 70$, $E = 3V/2 = 105$ bonds, $F = 2 - V + E = 37$
faces, so $37 - 12 = 25$ hexagons.
\end{solution}

\begin{solution}{exo:b3:inorganic-materials:5}
$\sqrt2\,(r_B + r_O) = 1.4142 \times \qty{200.5}{pm} = \qty{283.5}{pm}$. $\ce{SrTiO3}$: $284/283.5 = 1.00$, the ideal cubic perovskite.
$\ce{BaTiO3}$: $301/283.5 = 1.06$, titanium too small for its octahedron, off-centre: \omterm{def:b3:inorganic-materials:ferroelectric}{ferroelectric}.
$\ce{CaTiO3}$: $274/283.5 = 0.97$, calcium too small for its cavity: tilted octahedra, lower
symmetry.
\end{solution}

\begin{solution}{exo:b3:inorganic-materials:6}
Si/Al $= 106/86 = 1.23$. $M = 86 \times 23.0 + 86 \times 27.0 + 106 \times 28.1 + 384 \times 16.0 = \qty{13422.6}{g/mol}$; 43 $\ce{Ca^2+}$ per
formula: $43/\qty{13422.6}{g/mol} = \qty{3.20}{mmol/g}$.
\end{solution}

\begin{solution}{exo:b3:inorganic-materials:7}
Nearest-neighbour distance $d = a/\sqrt2 = \qty{288.4}{pm}$. A cluster of $n$ shells is $(2n + 1)d$ across:
$2n + 1 \approx 5/0.2884 = 17.3$, so $n = 8$ (\qty{4.9}{nm}). $N(8) = 2057$ atoms, $10 \times 64 + 2 = 642$ on the surface:
31\,\%.
\end{solution}

\begin{solution}{exo:b3:inorganic-materials:8}
$h^2/(8\mu R^2)$ with $\mu = \qty{9.109e-32}{kg}$: \qty{0.94}{eV} at \qty{2.0}{nm} and \qty{0.42}{eV} at \qty{3.0}{nm}. Emission at
$1.74 + 0.94 = \qty{2.68}{eV}$, \qty{463}{nm} (blue), and $1.74 + 0.42 = \qty{2.16}{eV}$, \qty{574}{nm} (yellow): the smaller dot
emits the bluer light.
\end{solution}

\begin{solution}{exo:b3:inorganic-materials:9}
$(10, 10)$: armchair, $n - m = 0$, metallic. $(10, 0)$: zigzag, 10 not a multiple of 3, semiconducting.
$(9, 0)$: zigzag, metallic. $(12, 8)$: neither (chiral), $n - m = 4$, semiconducting.
\end{solution}

\begin{solution}{exo:b3:inorganic-materials:10}
$M = 4 \times 65.4 + 13 \times 16.0 + 24 \times 12.0 + 12 \times 1.0 = \qty{769.6}{g/mol}$; $a^3 = (\qty{25.82e-8}{cm})^3 = \qty{1.721e-20}{cm^3}$;
$\rho = 8 \times 769.6/(\num{6.022e23} \times \num{1.721e-20}) = \qty{0.594}{g.cm^{-3}}$. A pitch is \qty{7140}{m^2}: one gram holds
$3000/7140 = 0.42$ pitch, and a teaspoon of a few grams more than a whole pitch.
\end{solution}

\begin{solution}{exo:b3:inorganic-materials:11}
$\dfrac{10n^2 + 2}{(10n^3 + 15n^2 + 11n + 3)/3} = \dfrac{30n^2 + 6}{10n^3 + 15n^2 + 11n + 3} \to \dfrac{30n^2}{10n^3} = \dfrac3n$. For $n = 8$: exact
$642/2057 = 0.31$, against $3/8 = 0.375$; the approximation overestimates small
clusters, whose inner shells are not negligible.
\end{solution}

\begin{solution}{exo:b3:inorganic-materials:12}
$2E = 3V$, $2E = 5p + 6h + 7s$, $F = p + h + s$. Euler: $V - E + F = 2$, with $V = 2E/3$: $F - E/3 = 2$, so
$p + h + s - (5p + 6h + 7s)/6 = 2$, that is $(p - s)/6 = 2$: $p - s = 12$. Each heptagon needs one extra
pentagon.
\end{solution}

\begin{solution}{pb:b3:inorganic-materials:1}
\textbf{1.} $\mathrm{Na_{12}Al_{12}Si_{12}O_{48}}$: $12 \times 23.0 + 12 \times 27.0 + 12 \times 28.1 + 48 \times 16.0 = \qty{1705.2}{g/mol}$.

\textbf{2.} $1705.2 + 27 \times 18.0 = \qty{2191.2}{g/mol}$.

\textbf{3.} Si/Al $= 12/12 = 1$.

\textbf{4.} No Al--O--Al links; with Si/Al $= 1$ every Si has four Al neighbours and every
Al four Si: strict alternation.

\textbf{5.} $-12$, one per aluminium.

\textbf{6.} $486.0/2191.2 = 22.2$\,\%.

\textbf{7.} $\ce{2Na+}(\text{zeolite}) + \ce{Ca^2+}(\text{aq}) \rightleftharpoons \ce{Ca^2+}(\text{zeolite}) + \ce{2Na+}(\text{aq})$.

\textbf{8.} Six (twelve charges).

\textbf{9.} $6/\qty{1705.2}{g/mol} = \qty{3.52}{mmol/g}$.

\textbf{10.} $6/\qty{2191.2}{g/mol} = \qty{2.74}{mmol/g}$.

\textbf{11.} $\qty{3.52}{mmol/g} \times \qty{100.1}{g/mol} = \qty{352}{mg}$ of $\ce{CaCO3}$ per gram.

\textbf{12.} $\qty{2.0}{mmol/L} \times \qty{15}{L} = \qty{30}{mmol}$.

\textbf{13.} $\qty{30}{mmol}/\qty{2.74}{mmol/g} = \qty{11}{g}$.

\textbf{14.} $\qty{10.9}{g}/0.60 = \qty{18}{g}$.

\textbf{15.} $\qty{60}{mmol}$ of $\ce{Na+}$, $\qty{0.060}{mol} \times \qty{23.0}{g/mol} = \qty{1.4}{g}$.

\textbf{16.} Phosphates in waste water feed algae in rivers and lakes; the blooms
and their decay use up the dissolved oxygen (eutrophication).

\textbf{17.} The smaller $\ce{Mg^2+}$ holds its water molecules more tightly; it must shed them
to pass the window and reach the sites, which is slow at washing temperatures.

\textbf{18.} The bare ion is about \qty{0.20}{nm} across, half the window.

\textbf{19.} The hydrated ion is larger than the window: it must shed part of its
water shell to pass, and the oxygens of the framework take the place of the
water.

\textbf{20.} The framework is negatively charged and exchanges only cations; a
large anion with a long chain is repelled and cannot pass the window, so it stays
in the water, where it does its work.

\textbf{21.} Dehydrated, the \omterm{def:b3:inorganic-materials:zeolite}{zeolite} takes up water strongly into its cages; only
molecules smaller than the window enter, so it separates molecules by size:
a sieve on the molecular scale.

\textbf{22.} The larger $\ce{K+}$ ions sit in the windows and narrow them: only smaller
molecules, such as water, are admitted.

\textbf{23.} The slim, linear \textit{para} isomer fits the channels and diffuses fast; the
wider \textit{ortho} and \textit{meta} isomers diffuse slowly, stay and isomerise.

\textbf{24.} Anhydrous \omterm{def:b3:inorganic-materials:zeolite}{zeolite} A can take up, in theory, \qty{3.52}{mmol} of $\ce{Ca^2+}$ per gram.
\end{solution}
